\documentclass[11pt]{article}
\usepackage[margin=1in]{geometry}
\usepackage{amsmath,amssymb,amsthm}
\usepackage[hidelinks]{hyperref}

\newtheorem{theorem}{Theorem}
\newtheorem{lemma}[theorem]{Lemma}
\newtheorem{corollary}[theorem]{Corollary}
\theoremstyle{remark}
\newtheorem{remark}[theorem]{Remark}

\newcommand{\Inf}{\mathrm{Inf}}
\newcommand{\Var}{\mathrm{Var}}
\newcommand{\E}{\mathbb{E}}
\newcommand{\poly}{\mathrm{poly}}
\newcommand{\polylog}{\mathrm{polylog}}
\newcommand{\PromiseBQP}{\mathsf{PromiseBQP}}
\newcommand{\PromiseBPP}{\mathsf{PromiseBPP}}

\title{On the Complexity of Finding an Influential Variable}
\author{Lance Fortnow}
\date{November 10, 2026}

\begin{document}
\maketitle

\begin{abstract}
The Aaronson--Ambainis conjecture, now a theorem, says every bounded degree-$d$ polynomial
$p:\{0,1\}^n\to[0,1]$ has a variable of influence $\poly(\Var(p)/d)$. The proof is
existential. We ask how hard it is to \emph{find} such a variable. When $p$ is given as
the acceptance probability of a quantum circuit, the search problem is
$\PromiseBQP$-complete: easy for a quantum computer and $\PromiseBQP$-hard even for
degree-one polynomials. When $p$ is given as a black box, a variable of influence
$\poly(\Var(p)/d)$ can be found with $\poly(d/\Var(p))\cdot\log n\log\log n$ queries, and
$\Omega(\log n)$ queries are necessary. All of the difficulty lies in evaluating $p$,
not in the search.
\end{abstract}

\section{Background}

Write $p(x)=\sum_{U\subseteq[n]}\hat p(U)\chi_U(x)$ in the $\pm1$ Fourier basis, so that
\[
\Inf_i(p)=\sum_{U\ni i}\hat p(U)^2=\tfrac14\,\E_x\big[(p(x)-p(x^{\oplus i}))^2\big],
\qquad
\Var(p)=\sum_{U\neq\emptyset}\hat p(U)^2 .
\]

Aaronson and Ambainis~\cite{AA14} conjectured that every polynomial
$p:\{0,1\}^n\to[0,1]$ of degree~$d$ satisfies
\[
\max_i\Inf_i(p)\ge\poly(\Var(p)/d).
\]
Turturean~\cite{Tur26} recently proved this. Blanc, Docter, Strassle, and
Tan~\cite{BDST-exp} give an exposition that obtains the bound
$\Omega(\Var(p)^8/d^{12})$; the original proof gives $\Omega(\Var(p)^{5/2}/d^6)$.

The motivation is the simulation of quantum query algorithms. By Beals et
al.~\cite{BBCMW01}, the acceptance probability of a $t$-query quantum algorithm is a
polynomial of degree at most $2t$ in the input bits. A classical simulator can query the
most influential variable, restrict, and recurse until the polynomial is nearly
constant. The influence bound controls the number of \emph{queries} this simulator
makes. The exposition does not discuss the cost of \emph{finding} the influential
variable. That cost appears only implicitly, through the complexity-class consequences:
Aaronson and Ambainis note that a random-oracle separation of $\mathsf{BQP}$ from
$\mathsf{HeurP}$ would require $\mathsf{P}\neq\mathsf{P}^{\#\mathsf{P}}$, and Blanc,
Docter, Strassle, and Tan~\cite{BDST26} strengthen this to
$\PromiseBPP\neq\PromiseBQP$. This bears on the question of Fortnow and
Rogers~\cite{FR99} of whether $\mathsf{BPP}\neq\mathsf{BQP}$ relative to a random oracle.

We consider two ways $p$ might be presented.

\section{When \texorpdfstring{$p$}{p} is given by a quantum circuit}

Suppose we are given a quantum circuit $Q$ that makes queries to an input
$x\in\{0,1\}^n$, and let $p(x)=\Pr[Q^x\text{ accepts}]$.

\begin{theorem}[Upper bound]\label{thm:circuit-upper}
There is a quantum algorithm that, given $Q$ and $\varepsilon>0$, runs in time
$\poly(|Q|,n,1/\varepsilon)$ and with high probability outputs $i$ with
$\Inf_i(p)\ge\max_j\Inf_j(p)-\varepsilon$.
\end{theorem}

\begin{proof}
Since $x^{\oplus i}$ is uniform when $x$ is,
\[
\Inf_i(p)=\tfrac12\Big(\E_x\big[p(x)^2\big]-\E_x\big[p(x)\,p(x^{\oplus i})\big]\Big).
\]
Running $Q$ twice independently on the same $x$, both runs accept with probability
$p(x)^2$; running it on $x$ and on $x^{\oplus i}$, both accept with probability
$p(x)p(x^{\oplus i})$. So each term is the acceptance probability of a small quantum
circuit and can be estimated to additive $\varepsilon/4$ with $O(\varepsilon^{-2}\log n)$
runs. Estimate all $n$ influences and output the largest.
\end{proof}

$\Var(p)$ can be estimated the same way, and restricting a variable amounts to hardwiring
it into $Q$. Hence every step of the Aaronson--Ambainis simulator, including the test of
whether $p$ is close to constant, is a $\PromiseBQP$ computation.

\begin{theorem}[Lower bound]\label{thm:circuit-lower}
Let $0<c<1/4$ be a constant. Finding a variable of influence at least $c$, given a quantum
circuit computing a degree-$1$ polynomial on two variables with $\Var(p)\ge1/8$, is
$\PromiseBQP$-hard.
\end{theorem}

\begin{proof}
Let $C$ be a $\PromiseBQP$ circuit, amplified so that its acceptance probability $r$
satisfies $r\ge1-2^{-k}$ or $r\le2^{-k}$, with $k$ large enough that $2^{-2k}/4<c\le(1-2^{-k})^2/4$.
Let $Q^x$ run $C$ and output $x_1$ if $C$ accepts and $x_2$ otherwise. Then
$p(x)=r\,x_1+(1-r)\,x_2$, which has degree~$1$,
\[
\Inf_1(p)=\tfrac{r^2}{4},\qquad \Inf_2(p)=\tfrac{(1-r)^2}{4},\qquad
\Var(p)=\tfrac{r^2+(1-r)^2}{4}\ge\tfrac18 .
\]
In the yes case only $x_1$ has influence at least $c$; in the no case only $x_2$ does.
The variable returned decides $C$.
\end{proof}

Since $d=1$ and $\Var(p)\ge1/8$, the target $\poly(\Var(p)/d)$ in the
Aaronson--Ambainis bound is a constant, so Theorem~\ref{thm:circuit-lower} applies to
exactly the guarantee the simulator needs. Finding the \emph{most} influential variable
is hard by the same reduction.

\begin{remark}
The circuit $C$ in the reduction never looks at $x$; the hardness is an unrelativized
$\mathsf{BQP}$ computation planted inside the polynomial. This is the ``easy direction''
(an unrelativized separation gives a random-oracle separation) in miniature, and the
Aaronson--Ambainis theorem says it is the only obstacle: with a $\PromiseBQP$ oracle the
classical simulator runs in polynomial time. The worst-case question is therefore
settled. What remains open is whether finding becomes easier when hardness cannot be
planted, for instance for circuits whose non-query gates are Clifford or of constant
depth, or on average over natural distributions of circuits. The existence proof gives
no guidance, since none of its steps is constructive.
\end{remark}

\section{When \texorpdfstring{$p$}{p} is given as a black box}

Now suppose we may query $p(x)$ for any $x$ we choose. Everything below also works in the
weaker model where a query returns a coin of bias $p(x)$, provided the same point can be
queried twice.

\begin{theorem}\label{thm:blackbox}
Let $\tau=\poly(\Var(p)/d)$ be the Aaronson--Ambainis bound. There is a randomized
algorithm that, with
$O\big(d\,\Var(p)\,\tau^{-3}\log n\cdot\log(d\log n/\tau)\big)=\poly(d/\Var(p))\cdot
\log n\cdot\log\log n$ queries to $p$, outputs with high probability a variable $i$ with
$\Inf_i(p)\ge\tau/2$. Conversely, $\Omega(\log n)$ queries are necessary.
\end{theorem}

The lower bound is immediate: for $p=x_i$ with $i$ hidden, each query returns one bit
and $i$ is the only influential variable, so $\log_2 n$ queries are needed.

For the upper bound we use the influence of a set of variables, as in junta
testing~\cite{FKRSS04,Blais09},
\[
\Inf_S(p)=\sum_{U\cap S\neq\emptyset}\hat p(U)^2=\tfrac12\,\E\big[(p(x)-p(y))^2\big],
\]
where $x$ is uniform and $y$ is $x$ with the coordinates in $S$ rerandomized. Two
properties matter. It is monotone, so $\Inf_S(p)\ge\Inf_i(p)$ whenever $i\in S$. And
it is dominated by the influence mass $W(S)=\sum_{i\in S}\Inf_i(p)=\sum_U|U\cap
S|\,\hat p(U)^2$, which is additive over disjoint sets and, because $p$ has degree $d$,
has total
\[
W([n])=\sum_U|U|\,\hat p(U)^2\le d\cdot\Var(p).
\]
Since $(p(x)-p(y))^2\in[0,1]$, $\Inf_S(p)$ can be estimated to additive $\varepsilon$ with
confidence $1-\gamma$ using $O(\varepsilon^{-2}\log(1/\gamma))$ queries, independent of
$|S|$. With coin access, write $\Inf_S(p)=\E[p(x)^2]-\E[p(x)p(y)]$ and estimate each
term from pairs of independent coins.

\begin{proof}[Proof of Theorem~\ref{thm:blackbox}]
Run the branch-and-prune search of Goldreich--Levin~\cite{GL89} and
Kushilevitz--Mansour~\cite{KM93}. Start with the active list $\{[n]\}$. At each of
$\lceil\log_2 n\rceil$ levels, split every active set in half, estimate $\Inf_S(p)$ for
each half to additive $\tau/4$, and keep the halves whose estimate exceeds $\tau/2$. At
the end, output any surviving singleton.

Assume all estimates are accurate. If $\Inf_{i^*}(p)\ge\tau$, every set containing
$i^*$ has $\Inf_S(p)\ge\tau$ by monotonicity and survives, so $\{i^*\}$ is reached. Every
surviving set has $\Inf_S(p)>\tau/4$, and in particular every surviving singleton has
influence above $\tau/4$; re-estimating the survivors to additive $\tau/8$ and outputting
the best one gives influence at least $\tau/2$. The active sets at a level are disjoint,
so their number is at most
\[
\frac{\sum_S \Inf_S(p)}{\tau/4}\le\frac{\sum_S W(S)}{\tau/4}\le\frac{4d\,\Var(p)}{\tau}.
\]
That gives $O(d\,\Var(p)\,\tau^{-1}\log n)$ estimates, each with
$O(\tau^{-2}\log(d\log n/\tau))$ queries to make a union bound work. (The algorithm needs
$d$ and a lower bound on $\Var(p)$ to set $\tau$; the variance can itself be estimated
with $\poly(1/\Var(p))$ queries.)
\end{proof}

The degree bound is used only to bound the total mass $W([n])$, and hence the width of
the search. This is the same argument Blanc, Docter, Strassle, and Tan~\cite[Prop.~10.6]{BDST26}
use to find a variable of large \emph{query weight}, with sums of query weights in
place of set influences.

\begin{remark}[Influences relative to a random oracle]
In~\cite{BDST26} the authors remark that the Aaronson--Ambainis subroutine for finding
a variable of high influence is in $\mathsf{NP}^{\#\mathsf{P}}$, whereas their
subroutine for query weights is in $\PromiseBQP$, and they cite this as an algorithmic
advantage of query weights. Combining the two-copy estimator in the proof of
Theorem~\ref{thm:circuit-upper} with the search above appears to put the influence
version in $\PromiseBQP$ as well. When the input is a random oracle
$O\in\{0,1\}^N$ and $Q$ makes $t$ queries, $\Inf_S$ for a dyadic interval $S$ is
$\E[p(O)^2]-\E[p(O)p(O')]$, a difference of average acceptance probabilities of
two-copy circuits. The random oracles $O$ and $O'$ can be simulated with
$2t$-wise independent functions~\cite{Zha12}, and the search needs only
$\poly(t)\cdot\log N$ such estimates. So the Aaronson--Ambainis regularity lemma also appears
to be algorithmic under $\PromiseBPP=\PromiseBQP$. If correct, the advantage of query
weights lies in their easier existence proof, not in their algorithmic tractability.
\end{remark}

\section{Related work}

Aaronson and Ambainis~\cite{AA14} made their regularity lemma algorithmic under
$\mathsf{P}=\mathsf{P}^{\#\mathsf{P}}$, finding influential variables by computing
influences exactly. Blanc, Docter, Strassle, and Tan~\cite{BDST26} replace influences
by query weights, a property of the algorithm rather than of $p$. They find heavy
variables with a $\PromiseBQP$ subset-sum estimator and branch-and-prune, which gives
logarithmic dependence on $N$, and conclude that a random-oracle separation
$\PromiseBQP^O\not\subseteq\PromiseHeurBPP^O$ would imply
$\PromiseBPP\neq\PromiseBQP$. None of these works considers hardness of finding.

Set influences and their estimation with $O(\varepsilon^{-2})$ queries, independent
of $|S|$, go back to Fischer, Kindler, Ron, Safra, and Samorodnitsky~\cite{FKRSS04} and
Blais~\cite{Blais09} in junta testing. Blais, Canonne, Eden, Levi, and
Ron~\cite{BCELR18} introduce the influence of a random $\rho$-biased subset of a set,
the same noise statistic one gets by flipping coordinates of $S$ with small
probability.

For learning rather than search, Eskenazis and Ivanisvili~\cite{EI22} show that
bounded degree-$d$ functions can be learned to $L_2$ error $\varepsilon$ from
$\log(n/\delta)\,\varepsilon^{-d-1}C^{d^{3/2}\sqrt{\log d}}$ uniformly random
queries, with the $\log n$ dependence optimal. Since $|\sqrt{\Inf_i(h)}-\sqrt{\Inf_i(p)}|\le\|h-p\|_2$,
this already finds an influential variable nonadaptively with $O(\log n)$ queries, but
with exponential dependence on $d$; Theorem~\ref{thm:blackbox} trades nonadaptivity
for polynomial dependence on $d$.

On the quantum side, Fourier sampling via the Bernstein--Vazirani circuit estimates
the influences of a Boolean function~\cite{FAH13,LY15}: position $i$ of the sample is
$1$ with probability $\Inf_i(f)$.

\section{Conclusion}

Combining the two settings, searching for an influential variable costs only
$O(\log n\log\log n)\cdot\poly(d/\Var(p))$ evaluations of $p$, while evaluating $p$ from a
quantum circuit is $\PromiseBQP$-hard even in the simplest case. The barrier in the
Blanc--Docter--Strassle--Tan result sits at $\PromiseBPP$ versus $\PromiseBQP$ because
of evaluation, not search.

\section*{AI disclosure}

This note was produced in a single conversation with Claude Opus~5.5, a large language
model developed by Anthropic, on November 10, 2026. The author posed the questions: whether
the Blanc--Docter--Strassle--Tan exposition addresses the complexity of finding an
influential variable, and how hard that problem is when $p$ is given by a quantum circuit
and when it is given as a black box. Claude proposed the theorems and proofs in
Sections~2 and~3, carried out the literature search behind Section~4 (including the
observation in Remark~5), and wrote the text and \LaTeX. The literature search led to the
simpler branch-and-prune proof of Theorem~\ref{thm:blackbox}, which replaced a more
complicated argument in an earlier draft. The arguments have not been independently
verified, and some bibliographic details were supplied from the model's own knowledge
rather than checked against the sources.

\begin{thebibliography}{BBCMW01}

\bibitem[AA14]{AA14}
Scott Aaronson and Andris Ambainis.
\newblock The need for structure in quantum speedups.
\newblock {\em Theory of Computing}, 10(6):133--166, 2014.

\bibitem[BBCMW01]{BBCMW01}
Robert Beals, Harry Buhrman, Richard Cleve, Michele Mosca, and Ronald de Wolf.
\newblock Quantum lower bounds by polynomials.
\newblock {\em Journal of the ACM}, 48(4):778--797, 2001.

\bibitem[BCELR18]{BCELR18}
Eric Blais, Cl\'ement~L. Canonne, Talya Eden, Amit Levi, and Dana Ron.
\newblock Tolerant junta testing and the connection to submodular optimization and
  function isomorphism.
\newblock In {\em Proceedings of the 29th ACM-SIAM Symposium on Discrete Algorithms
  (SODA)}, 2018. arXiv:1607.03938.

\bibitem[BDST26a]{BDST26}
Guy Blanc, Jordan Docter, Carmen Strassle, and Li-Yang Tan.
\newblock Quantum speedups require structure or depth.
\newblock In {\em Proceedings of the 67th IEEE Symposium on Foundations of Computer
  Science (FOCS)}, 2026. arXiv:2608.19158.

\bibitem[BDST26b]{BDST-exp}
Guy Blanc, Jordan Docter, Carmen Strassle, and Li-Yang Tan.
\newblock Proof of the {Aaronson--Ambainis} conjecture: An exposition.
\newblock Manuscript, November 2026.
\newblock \url{https://guyblanc.github.io/AA_Exposition.pdf}.

\bibitem[Bla09]{Blais09}
Eric Blais.
\newblock Testing juntas nearly optimally.
\newblock In {\em Proceedings of the 41st ACM Symposium on Theory of Computing
  (STOC)}, pages 151--158, 2009.

\bibitem[EI22]{EI22}
Alexandros Eskenazis and Paata Ivanisvili.
\newblock Learning low-degree functions from a logarithmic number of random queries.
\newblock In {\em Proceedings of the 54th ACM Symposium on Theory of Computing
  (STOC)}, pages 203--207, 2022.

\bibitem[FAH13]{FAH13}
Dominik~F. Floess, Erika Andersson, and Mark Hillery.
\newblock Quantum algorithms for testing and learning {Boolean} functions.
\newblock {\em Mathematical Structures in Computer Science}, 23(2):386--398, 2013.

\bibitem[FKRSS04]{FKRSS04}
Eldar Fischer, Guy Kindler, Dana Ron, Shmuel Safra, and Alex Samorodnitsky.
\newblock Testing juntas.
\newblock {\em Journal of Computer and System Sciences}, 68(4):753--787, 2004.

\bibitem[FR99]{FR99}
Lance Fortnow and John Rogers.
\newblock Complexity limitations on quantum computation.
\newblock {\em Journal of Computer and System Sciences}, 59(2):240--252, 1999.

\bibitem[GL89]{GL89}
Oded Goldreich and Leonid Levin.
\newblock A hard-core predicate for all one-way functions.
\newblock In {\em Proceedings of the 21st ACM Symposium on Theory of Computing}, pages
  25--32, 1989.

\bibitem[KM93]{KM93}
Eyal Kushilevitz and Yishay Mansour.
\newblock Learning decision trees using the {Fourier} spectrum.
\newblock {\em SIAM Journal on Computing}, 22(6):1331--1348, 1993.

\bibitem[LY15]{LY15}
Hongwei Li and Li Yang.
\newblock A quantum algorithm for approximating the influences of {Boolean}
  functions and its applications.
\newblock {\em Quantum Information Processing}, 14:1787--1797, 2015.

\bibitem[Tur26]{Tur26}
David Turturean.
\newblock A proof of the {Aaronson--Ambainis} conjecture.
\newblock Manuscript, November 2026.
\newblock \url{https://psi.inc/research/aaronson-ambainis/}.

\bibitem[Zha12]{Zha12}
Mark Zhandry.
\newblock Secure identity-based encryption in the quantum random oracle model.
\newblock In {\em Advances in Cryptology -- CRYPTO 2012}, pages 758--775, 2012.

\end{thebibliography}

\end{document}
